\documentclass[11pt]{article}
\usepackage[margin=1.05in]{geometry}
\usepackage{amsmath,amssymb,amsthm}
\usepackage{booktabs}
\usepackage{array}
\usepackage{microtype}
\usepackage{xcolor}
\usepackage[hidelinks]{hyperref}
\usepackage{url}
\usepackage{enumitem}

\newtheorem{theorem}{Theorem}
\newtheorem{lemma}[theorem]{Lemma}
\newtheorem{corollary}[theorem]{Corollary}
\newtheorem{proposition}[theorem]{Proposition}
\newtheorem{conjecture}[theorem]{Conjecture}
\theoremstyle{definition}
\newtheorem{definition}[theorem]{Definition}
\theoremstyle{remark}
\newtheorem{remark}[theorem]{Remark}

\newcommand{\Phichain}{\Phi}
\newcommand{\nfour}{n_4}
\newcommand{\ntwo}{n_2}
\newcommand{\mass}{\mathrm{mass}}
\newcommand{\Str}{\mathrm{S}}

\title{Perfect 2048:\\ the fewest moves, the highest score,\\ and what the undo button buys}
\author{Dom the Developer\\ \small \url{https://github.com/DomTheDeveloper/2048-undo}}
\date{Draft, September 2026}

\begin{document}
\maketitle

\begin{abstract}
We study 2048 under \emph{possible play}: the spawn sequence is treated as
part of the plan, which is exactly what a player with a re-randomising undo
button can realise, and exactly what ``possible with positive probability''
means for the ordinary game. Three results. (1)~On any board, the tile
$2^n$ cannot appear before move $(2^n-8)/4+(n-2)$; the bound is attained,
giving $519$ for the 2048 tile (recovering Lees-Miller's Markov-chain
figure) and $32{,}781$ for the largest tile $131072$ on the $4\times4$
board, and $65{,}533$ moves for the complete descending chain
$131072,65536,\dots,4$. (2)~On any board with $c$ cells the score is at
most $\Phi_c-4c$ where $\Phi_c=\sum_{k=2}^{c+1}(k-1)2^k$: reaching the
top position needs at least one spawned~4 for every one of its $c$ tiles,
and each spawned~4 forfeits exactly four points. The proof is a
register-allocation (pebbling) argument. For $4\times4$ this is
$3{,}932{,}100$. Exhaustive enumeration shows the bound is tight on every
board with at most nine cells---$2\times2$, $2\times3$, $2\times4$,
$3\times3$---and reproduces, to the digit, the $48{,}713{,}519$ positions of
$3\times3$ reported by Yamashita, Kaneko and Nakayashiki. For $4\times4$ we
give an explicit line scoring $3{,}925{,}224$ in $129{,}333$ moves; whether
the remaining $6{,}876$ points can be collected is open, and we conjecture
they can. (3)~The perfect lines are shipped as data and replayed in a
browser; with honest 90/10 spawns the undo button must be pressed about
$1.94$ million times to play the $32{,}781$-move line, a number we compute
exactly. We also strongly solve the boards up to nine cells (optimal
expected score, tile probabilities) and benchmark an expectimax player
against the classical heuristics on the full board.
\end{abstract}

\section{Introduction}

2048 \cite{cirulli2014} is played on a $4\times4$ board. A move slides all
tiles in one of four directions; equal tiles that collide merge into their
sum, each tile merging at most once per move; then one new tile appears in
a uniformly random empty cell, a 2 with probability $0.9$ and a 4 with
probability $0.1$. The score is the sum of all merged values. The game
ends when no move changes the board.

Alok Menghrajani's fork \cite{menghrajani2014} adds an undo button that
re-randomises the spawn on replay. That small change turns the game into a
different object: a player who is willing to press undo often enough can
realise \emph{any} spawn sequence, and so can play any game that is possible
in principle. This paper is about those games. Three questions organise it:
\begin{enumerate}[nosep]
\item What is the fewest moves in which the largest tile can be built?
\item What is the highest score any game can reach?
\item What does the undo button cost, in presses, to play such a game?
\end{enumerate}

The first two questions have folklore answers. For the move count, the
observation that every move adds a spawned 2 or 4 to the board gives a
lower bound immediately; the cascade that finishes the tile adds a few more
moves, and it was not obvious that the sum is exact. For the score, several
numbers circulate: $3{,}932{,}164$ (pure mass accounting), $3{,}932{,}156$
(two ``structurally forced'' 4-spawns), and $3{,}932{,}100$ (one forced
4-spawn per tile of the final position). Our own earlier README argued
for a ceiling of $3{,}932{,}156$ and claimed geometry prevents paying it.
That claim was half right: geometry does prevent it, but the right ceiling
is $3{,}932{,}100$, and we prove it.

\paragraph{Contributions.}
\begin{itemize}[nosep]
\item Theorem~\ref{thm:moves}: a complete proof that $2^n$ needs at least
$(2^n-8)/4+(n-2)$ moves on any board, with equality for $n=11$ (519) and
$n=17$ ($32{,}781$) on $4\times4$, and for every tile on every board of at
most nine cells.
\item Theorem~\ref{thm:score}: on any $c$-cell board, score $\le\Phi_c-4c$.
The argument identifies a merge tree's cell requirement with its Strahler
number (the Ershov register count) and shows that the $c$ final tiles must
pay their tight moments in a forced order.
\item Exact values (Section~\ref{sec:exact}): the bound is attained on
$2\times2$, $2\times3$, $2\times4$ and $3\times3$; the enumerator agrees
with the published $3\times3$ state count.
\item Section~\ref{sec:fourbyfour}: for $4\times4$, an explicit
$129{,}333$-move line scoring $3{,}925{,}224$ with $1{,}735$ spawned 4s,
where the theorem allows $16$; we locate where the construction spends its
4s and conjecture the ceiling is attainable.
\item Section~\ref{sec:undo}: the exact expected number of undo presses
for each shipped line under honest odds, and the opening problem.
\item Section~\ref{sec:honest}: a strong solution of the small boards
(expected score, tile probabilities) and an expectimax player for the
full board measured against the classical heuristics and an adversarial
tile generator.
\end{itemize}

\section{Related work}

\paragraph{Complexity and combinatorics.} Langerman and Uno
\cite{langerman2016} showed that deciding whether a given position can be
played to a given tile is NP-hard on $m\times n$ boards for Threes, 1024
and 2048; Abdelkader, Acharya and Dasler \cite{abdelkader2016} showed the
variant without new tiles is still hard; Mehta \cite{mehta2014} gave a
PSPACE-completeness result for a related decision problem. Eppstein
\cite{eppstein2018} analysed the maximum score in a relaxed model of
2048-like games in which any set of cells whose values sum to a tile value
may be merged in one step, proving a min--max theorem that equates the
maximum score with the smallest value for which greedy change-making needs
a given number of coins; the binary-counter strategy every strong 2048
player follows is, in that light, greedy change-making in base~2. The
relaxed model has only the small spawn value. The real game's 4-spawns are
precisely the resource that buys an extra cell, and Theorem~\ref{thm:score}
prices that resource: four points each, and exactly $c$ of them on the way
to the top position.

\paragraph{Learning agents.} Szubert and Ja\'skowski \cite{szubert2014}
introduced temporal-difference learning of $n$-tuple networks for 2048; Wu
et al.\ \cite{wu2014} added multi-stage learning; Ja\'skowski
\cite{jaskowski2018} reached the strongest reported agents; Matsuzaki and
co-authors studied the systematic selection of $n$-tuples and
neural-network players \cite{oka2016,matsuzaki2016,matsuzaki2017,matsuzaki2019};
Guei, Wei and Wu \cite{guei2020,guei2021} used 2048-like games for
teaching reinforcement learning and proposed optimistic TD learning.
Slizkov \cite{slizkov2023} proved rigorous bounds on honest play: a
strategy reaching 2048 with probability at least $0.99969$, and 256 with
certainty.

\paragraph{Exact solving.} Lees-Miller \cite{leesmiller2017,leesmiller2017b,leesmiller2018}
analysed the game with Markov chains and Markov decision processes,
computed the minimum of 519 moves to the 2048 tile, and enumerated the
$2\times2$ and $3\times3$ games. Yamashita, Kaneko and Nakayashiki
\cite{yamashita2022} strongly solved $3\times3$ ($48{,}713{,}519$ positions)
and Kaneko and Yamashita \cite{kaneko2026} strongly solved $4\times3$
($1{,}152{,}817{,}492{,}752$ positions, $739{,}648{,}886{,}170$ afterstates,
optimal expected score about $50{,}724.26$ from a two-2s start), the key
technique being to partition the state space by the sum of tile values
(the ``age'' of a state). Our enumerator uses the same grading and, on
$3\times3$, reproduces their position count exactly
(Section~\ref{sec:exact}).

\section{Preliminaries}

A board is a finite set of $c$ cells with a slide structure (rows and
columns in the rectangular case; the arguments below use only the cell
count). Tiles carry values $2^k$, $k\ge1$; we call $k$ the \emph{rank}. A
\emph{play} is a sequence of moves, each consisting of a slide (which
performs its merges) followed by a \emph{spawn} of a 2 or a 4 into an empty
cell. The game starts with two spawns. Let $\ntwo$ and $\nfour$ be the
numbers of spawned 2s and 4s so far, counting the two starting tiles.

\begin{definition}[possible play]
In \emph{possible play} the spawn values and cells are chosen by the
player. Every finite spawn sequence has positive probability in the honest
game, so a position is reachable in possible play if and only if it is
reachable with positive probability in the honest game; and a player with
a re-randomising undo button can realise any possible play.
\end{definition}

For a position $B$ write $\mass(B)=\sum_{t\in B}2^{k_t}$ and
\[
\Phichain(B)=\sum_{t\in B}(k_t-1)\,2^{k_t}.
\]

\begin{lemma}[mass ledger]\label{lem:mass}
Slides conserve mass and every spawn adds its value:
$\mass(B)=2\ntwo+4\nfour$ at all times. After $m$ moves the mass is at
most $8+4m$, with equality if and only if every spawn so far, the two
starting tiles included, was a~4.
\end{lemma}

\begin{lemma}[score identity]\label{lem:score}
At every moment of every play, $\mathrm{score}=\Phichain(B)-4\nfour$.
\end{lemma}
\begin{proof}
By induction over events. Initially the score is $0$; a starting 2
contributes $0$ to $\Phichain$, a starting 4 contributes $4$ and raises
$4\nfour$ by $4$. A merge of two $2^k$ tiles into $2^{k+1}$ scores
$2^{k+1}$ and changes $\Phichain$ by $k2^{k+1}-2(k-1)2^k=2^{k+1}$. A spawned
2 changes nothing on either side; a spawned 4 adds $4$ to $\Phichain$ and
$4$ to $4\nfour$.
\end{proof}

So every spawned 4 costs exactly four points relative to what the final
position ``looks like it is worth'', and the maximum score of a game is
$\max_B\bigl(\Phichain(B)-4\,\nfour^{\min}(B)\bigr)$ over reachable
positions $B$, where $\nfour^{\min}(B)$ is the fewest spawned 4s on any play
reaching $B$.

\paragraph{Merge trees.} Every tile on the board is the root of a binary
\emph{merge tree}: its leaves are spawned tiles (value 2 or 4), each
internal node is a tile created by merging its two children. A tile of
rank $k$ has all leaves at depth $k-1$ (2-leaves) or $k-2$ (4-leaves). Two
facts about the game drive everything below:
\begin{itemize}[nosep]
\item \emph{Persistence.} A tile, once spawned or created, stays on the
board until it merges into its parent. Nothing else removes it.
\item \emph{One doubling per move.} A tile's value at most doubles in a
move, because each tile merges at most once per move.
\end{itemize}
Call the tiles present on the board that belong to the merge tree of a
future tile $T$ the \emph{pieces} of $T$. By persistence, from the first
spawn of a leaf of $T$ until $T$ is created, at least one piece of $T$ is
on the board; afterwards $T$ itself is.

\section{The fewest moves}

\begin{theorem}\label{thm:moves}
On any board and for any play, the tile $2^n$ ($n\ge3$) does not exist
before move
\[
M(n)=\frac{2^n-8}{4}+(n-2).
\]
Equality is attainable only if the game starts from two 4s and every spawn
outside the last $n-2$ moves is a~4.
\end{theorem}
\begin{proof}
Let $T$ be the first move after which a tile $2^n$ exists. A tile spawned
at move $T-j$ has value at most $4=2^2$ then, and by one doubling per move
its value at move $T$ is at most $2^{2+j}$; for it to be part of the $2^n$
we need $2+j\ge n$, i.e.\ $j\ge n-2$. Hence the $n-2$ tiles spawned at
moves $T-(n-3),\dots,T$ are not part of the $2^n$: they are still on the
board (persistence), as themselves or merged among each other, so the
position after move $T$ has mass at least $2^n$ plus the sum of those
$n-2$ spawn values.

Write $z$ for the number of spawned 2s among the $T$ move-spawns and $z'$
for the number among the last $n-2$ of them, so $z'\le z$. By
Lemma~\ref{lem:mass} the mass after move $T$ is $M_0+4T-2z$ with
$M_0\le8$ the starting mass, and the last $n-2$ spawns sum to $4(n-2)-2z'$.
Therefore
\[
M_0+4T-2z\;\ge\;2^n+4(n-2)-2z',
\qquad\text{so}\qquad
4T\;\ge\;2^n-8+4(n-2)+2(z-z')+(8-M_0)\;\ge\;2^n-8+4(n-2),
\]
which is the claim. Equality forces $M_0=8$ and $z=z'$.
\end{proof}

Attainability is a construction. For $n=11$, Lees-Miller's Markov chain
\cite{leesmiller2017} exhibits the $519=510+9$ move win on $4\times4$. For
$n=17$ we generated, verified and ship a $32{,}781$-move line
(Section~\ref{sec:lines}); Table~\ref{tab:exact} shows that the bound is
attained for every tile on every board up to nine cells.

\begin{corollary}\label{cor:chain}
On $4\times4$ the complete descending chain $131072,65536,\dots,4$ (one
power of two per cell, a dead position) needs exactly $65{,}533$ moves.
\end{corollary}
\begin{proof}
Its mass is $2^{18}-4=262{,}140$; by Lemma~\ref{lem:mass} at least
$(262{,}140-8)/4=65{,}533$ moves are needed and every spawn must be a~4.
The shipped line (Section~\ref{sec:lines}) does it in exactly that many.
Along it the $131072$ first forms at move $32{,}784$, three later than
$M(17)$: all-4 feeding needs a few consolidations that a 2-spawn during
the cascade would avoid, and the ledger absorbs them.
\end{proof}

\section{The highest score}

\subsection{A merge tree needs its Strahler number of cells}

The \emph{Strahler number} of a binary tree is $\Str(\text{leaf})=1$ and,
for a node with subtrees $A,B$, $\Str=\max(\Str A,\Str B)$ if they differ
and $\Str A+1$ if they are equal. It is the number of registers needed to
evaluate an expression tree without spilling (Ershov \cite{ershov1958},
Sethi--Ullman \cite{sethi1970}, Flajolet--Raoult--Vuillemin
\cite{flajolet1979}) and the black pebbling number of the tree
\cite{paterson1970}. Building a tile is evaluating its merge tree with
cells as registers, and the persistence of pieces makes the lower bound
short to prove directly.

\begin{lemma}\label{lem:strahler}
Let $T$ be a tile of rank $k\ge2$ with merge tree $\mathcal T$. Then
$\Str(\mathcal T)=k$ if every leaf is a 2, and $\Str(\mathcal T)=k-1$
if at least one leaf is a 4.
\end{lemma}
\begin{proof}
Induction on $k$. For $k=2$: a spawned 4 is a leaf ($\Str=1$), and
$2+2$ has $\Str=2$. For $k>2$ the children are trees of rank $k-1$ with
Strahler numbers in $\{k-2,k-1\}$, equal to $k-1$ exactly when all their
leaves are 2s. Two all-2 children give $\Str=k$; otherwise the maximum is
$k-1$ and, if both are $k-2$, $\Str=k-1$ as well.
\end{proof}

\begin{lemma}[tight moment]\label{lem:tight}
During the construction of a tile with merge tree $\mathcal T$ there is a
moment at which at least $\Str(\mathcal T)$ pieces of $\mathcal T$ are on
the board simultaneously.
\end{lemma}
\begin{proof}
Induction on $\mathcal T$; a leaf is its own piece. Let the root have
subtrees $A$ and $B$, roots created at times $t_A<t_B$ (if they are
created in the same move, either order works), merged at time $t^*>t_B$.
Let $\tau_A,\tau_B$ be tight moments of $A,B$ given by induction, so
$\tau_A\le t_A$ and $\tau_B\le t_B$. If $\Str A>\Str B$, or $\Str B>\Str A$,
the larger side's tight moment already shows $\max(\Str A,\Str B)$ pieces.
Suppose $\Str A=\Str B=s$; we need $s+1$.
\begin{itemize}[nosep]
\item If $\tau_B>t_A$: at $\tau_B$ the $s$ pieces of $B$ are on the board
and so is the finished root of $A$ (created at $t_A<\tau_B<t_B\le t^*$,
and persistent): $s+1$ pieces.
\item If $\tau_B\le t_A$ and $\tau_A>t_B$: symmetric, $s$ pieces of $A$
plus the root of $B$.
\item If $\tau_B\le t_A$ and $\tau_A\le t_B$: from $\tau_B$ until $t_B$
at least one piece of $B$ is on the board (persistence), and $\tau_A\le
t_A<t_B$; if $\tau_A\ge\tau_B$ then at $\tau_A$ there are $s$ pieces of $A$
and at least one of $B$. If instead $\tau_A<\tau_B\le t_A$, then at
$\tau_B$ there are $s$ pieces of $B$ and at least one piece of $A$
(present from $\tau_A$ until $t_A\ge\tau_B$).
\end{itemize}
In every case some moment shows $s+1$ pieces of $\mathcal T$.
\end{proof}

\subsection{The ceiling}

Let $\Phi_c=\sum_{k=2}^{c+1}(k-1)2^k=(c-1)2^{c+2}+4$, the value of
$\Phichain$ on the \emph{full chain} $2^{c+1},2^c,\dots,2^2$: one distinct
power of two per cell.

\begin{theorem}\label{thm:score}
On any board with $c$ cells, every play satisfies
$\mathrm{score}\le\Phi_c-4c$. Equality requires the final position to be
the full chain, reached with exactly one spawned 4 in the merge tree of
each of its $c$ tiles.
\end{theorem}
\begin{proof}
Fix a play and its final position with tiles $T_1,\dots,T_m$, $m\le c$, of
ranks $k_1,\dots,k_m$. Every spawned tile of the whole play is a leaf of
exactly one $T_i$, so $\nfour=\sum_i h_i'$ where $h_i'$ counts the 4-leaves
of $T_i$; put $h_i=\min(h_i',1)$. By Lemma~\ref{lem:strahler} the Strahler
number of $T_i$'s tree is $S_i=k_i-h_i$ (a tile of rank~1 is a spawned 2,
$S_i=1$).

Let $\tau_i$ be a tight moment of $T_i$ (Lemma~\ref{lem:tight}) and
relabel so that $\tau_1\le\tau_2\le\dots\le\tau_m$. At $\tau_p$ the board
holds $S_p$ pieces of $T_p$ and, for each $q<p$, at least one piece of
$T_q$ or $T_q$ itself (persistence, since $\tau_q\le\tau_p$). Hence
$S_p+(p-1)\le c$, i.e.
\[
k_p\;\le\;c+1-p+h_p\qquad(p=1,\dots,m).
\]
By Lemma~\ref{lem:score},
\[
\mathrm{score}=\sum_{p}(k_p-1)2^{k_p}-4\nfour\;\le\;\sum_{p}\Bigl[(k_p-1)2^{k_p}-4h_p\Bigr].
\]
For fixed $p$ the bracket, subject to $k_p\le c+1-p+h_p$, is largest with
$h_p=1$ and $k_p=c+2-p$, worth $(c+1-p)2^{c+2-p}-4$; the alternative
$h_p=0$, $k_p=c+1-p$ is worth $(c-p)2^{c+1-p}$, which is smaller for $p<c$
and equal ($=0$) for $p=c$. All the terms are nonnegative, so filling the
board ($m=c$) is best, and
\[
\mathrm{score}\;\le\;\sum_{p=1}^{c}\Bigl[(c+1-p)2^{c+2-p}-4\Bigr]
=\sum_{k=2}^{c+1}(k-1)2^k-4c=\Phi_c-4c.
\]
Equality forces $k_p=c+2-p$ and $h_p=1$ with $h'_p=1$ for every $p$.
\end{proof}

\begin{corollary}\label{cor:4x4}
On $4\times4$, $\mathrm{score}\le\Phi_{16}-64=3{,}932{,}164-64=3{,}932{,}100$.
The numbers $3{,}932{,}164$ and $3{,}932{,}156$ are not attainable.
\end{corollary}

\begin{remark}
The theorem uses only the number of cells, persistence, and one merge per
tile per move. It holds for any adjacency structure, any spawn rule that
offers the values 2 and 4, and in particular in possible play. The
constraint that actually binds is that a tile of rank $k$ built from 2s
alone needs $k$ cells at once, and the $c$ tiles of the top position must
take their tight moments in decreasing order of size on a board that the
larger ones are already sitting on.
\end{remark}

\subsection{Exact values on small boards}\label{sec:exact}

The mass of a position grows by exactly the spawned value each move, so
the state graph is a directed acyclic graph graded by mass, and can be
enumerated one layer at a time with only the two open layers in memory.
Kaneko and Yamashita \cite{kaneko2026} call this grading the age of a
state. We implemented it (\texttt{test/solve.js}) with positions reduced
by the board's symmetry group and, per position, the fewest spawned 4s
and the fewest moves on any play reaching it; the maximum score is then
$\max_B(\Phichain(B)-4\nfour^{\min}(B))$ by Lemma~\ref{lem:score}. The
same layered pass, run backwards, strongly solves the honest game: the
optimal expected score and, for each tile, the maximum probability of ever
building it.

\begin{table}[t]
\centering\small
\begin{tabular}{@{}lrrrrrr@{}}
\toprule
board & cells & positions & (raw) & max score & $\Phi_c-4c$ & $\Phi_c-8$ \\
\midrule
$2\times2$ & 4 & 110 & 662 & \textbf{180} & 180 & 188 \\
$2\times3$ & 6 & 21{,}752 & 85{,}844 & \textbf{1{,}260} & 1{,}260 & 1{,}276 \\
$2\times4$ & 8 & 4{,}980{,}767 & 19{,}920{,}382 & \textbf{7{,}140} & 7{,}140 & 7{,}164 \\
$3\times3$ & 9 & 48{,}713{,}519 & 388{,}921{,}077 & \textbf{16{,}352} & 16{,}352 & 16{,}380 \\
$4\times4$ & 16 & ? & ? & $\ge3{,}925{,}224$ & 3{,}932{,}100 & 3{,}932{,}156 \\
\bottomrule
\end{tabular}
\caption{Exact maximum scores in possible play. ``positions'' counts
reachable positions up to the board's symmetries (8 for squares, 4 for
rectangles), ``raw'' without reduction. On every board up to nine cells
the maximum equals the ceiling of Theorem~\ref{thm:score}, and the full
chain is reached with exactly $c$ spawned 4s. The $3\times3$ count agrees
with Yamashita et al.\ \cite{yamashita2022,kaneko2026}.}
\label{tab:exact}
\end{table}

Table~\ref{tab:exact} summarises. On every board tested the ceiling is
attained, always on the full chain, always with exactly one spawned 4 per
tile; the fewest moves to every tile $2^k$ with $k\ge3$ equals $M(k)$ of
Theorem~\ref{thm:moves}. The $3\times3$ enumeration gives
$48{,}713{,}519$ positions, the number Yamashita, Kaneko and Nakayashiki
report, which also fixes the convention of their count: positions up to
the eight symmetries of the square. (Lees-Miller's ``about 25 million''
for $3\times3$ counts differently.) The honest-game side of the same
computation is in Table~\ref{tab:honest}.

\begin{table}[t]
\centering\small
\begin{tabular}{@{}lrrl@{}}
\toprule
board & $\mathbb E[\text{score}]$, standard start & given a two-2s start & max probability of the largest tiles \\
\midrule
$2\times2$ & 66.96 & 67.70 & $16$: 96.25\%, \; $32$: 8.29\% \\
$2\times3$ & 480.26 & 480.99 & $64$: 93.93\%, \; $128$: 5.70\% \\
$2\times4$ & 2{,}641.87 & 2{,}642.60 & $256$: 86.98\%, \; $512$: 2.46\% \\
$3\times3$ & 5{,}468.49 & 5{,}469.21 & $512$: 73.68\%, \; $1024$: 1.13\% \\
\bottomrule
\end{tabular}
\caption{Strong solutions of the honest game (uniform empty cell, 90\%
twos). The expected score is under the score-maximising policy; each tile
probability is under the policy that maximises that probability. The
two-2s column uses Kaneko and Yamashita's convention (their $4\times3$
figure is $50{,}724.26$).}
\label{tab:honest}
\end{table}

\subsection{The $4\times4$ board}\label{sec:fourbyfour}

The $4\times4$ state space is out of reach of enumeration (Lees-Miller's
month-long run \cite{leesmiller2017b} enumerated $1.3\cdot10^{12}$ states
without finishing; the full game is far larger). What we have is a
constructive lower bound. A generator (\texttt{test/gen\_perfect.js})
walks a binary counter along the snake path into a corner, feeding 2s,
and grants a spawned 4 only when a bounded look-ahead shows every pure-2
continuation dying. It produced a line of $129{,}333$ moves ending on the
full chain with $\nfour=1{,}735$, hence a score of
\[
\Phi_{16}-4\cdot1{,}735=3{,}932{,}164-6{,}940=3{,}925{,}224,
\]
$99.825\%$ of the ceiling and $6{,}876$ points short of it. The identities
$\text{moves}=131{,}068-\nfour$ (all $129{,}333$ moves spawn 2 except the
4s; mass $262{,}140=8+2\cdot\text{moves}+2\nfour$) and
$\text{score}=3{,}932{,}164-4\nfour$ are verified by replay through the
game engine in all four corners and both spiral orientations.

Where do the $1{,}735$ go? The first act---building $131072$ on the
empty board---spends exactly one, at its tight moment. The whole gap is in
the second act, the chain $65536,\dots,4$ on the fifteen cells beside the
$131072$: $295$ spawned 4s while building the $65536$ and $1{,}439$ below
it, most of them with three to five empty cells still on the board. So the
generator is not being squeezed by the pebbling bound; it is being squeezed
by geometry within a cascade. When the counter overflows, the resulting
chain of merges takes one move per level, a tile spawns during each of
those moves, and the spawns land in the cells the cascade has just freed,
alongside the tail of the chain. The bounded look-ahead cannot always see
how to digest that junk with 2s alone and buys its way out with a 4.
Theorem~\ref{thm:score} says the price of the second act is at least
fifteen 4s; the small boards, which contain the same corner turns and the
same overflow cascades in miniature, all pay exactly the theorem's price.

\begin{conjecture}\label{conj}
The maximum score in $4\times4$ 2048 is $3{,}932{,}100$, attained by a
play with exactly sixteen spawned 4s.
\end{conjecture}

Two routes to settle it suggest themselves. The direct route is a better
search for the second act: the target is a schedule that, at each level
$k$ of the chain, builds the first copy of $2^{k-1}$ from 2s on the $k-1$
free cells (which Lemma~\ref{lem:strahler} permits) and the second copy
with a single 4 at its own tight moment, and that handles the overflow
junk explicitly. The indirect route is Kaneko and Yamashita's enumeration:
their layered pass over $4\times3$ can carry a two-byte ``fewest 4s''
field at negligible cost and would settle the $4\times3$ maximum, which
Theorem~\ref{thm:score} bounds by $\Phi_{12}-48=180{,}180$; we predict
equality.

\section{The undo button}\label{sec:undo}

\subsection{The shipped lines}\label{sec:lines}

Three lines were generated once and are shipped as data
(\texttt{js/perfect\_line.js}, two hex characters per move: direction,
spawn cell, spawn value): the $32{,}781$-move line to $131072$, the
$65{,}533$-move line to the full chain, and the $129{,}333$-move maximum
score line. The other three corners and the column-first spirals are the
seven symmetries of the square applied to the same data. A replay through
the real engine checks that every slide moves, every spawn cell is empty,
and the counts land exactly.

\subsection{What a perfect game costs in undos}

Under honest odds a spawn lands in a required cell with a required value
with probability $p=\tfrac{1}{e}\cdot q$, where $e$ is the number of empty
cells after the slide and $q\in\{0.9,0.1\}$ the value's probability. With
a re-randomising undo, the number of attempts is geometric and the expected
number of undos for the move is $1/p-1$; summing over the line gives the
exact expectation. The opening pair is the same problem once: two 4s in
two prescribed cells happen with probability $2\cdot\frac1{16}\cdot\frac1{15}\cdot0.01$,
one game in $12{,}000$, so the line to $131072$ starts with an expected
$11{,}999$ restarts.

\begin{center}\small
\begin{tabular}{@{}lrrrr@{}}
\toprule
line & moves & $\mathbb E[\text{undos in play}]$ & $\mathbb E[\text{opening restarts}]$ & total \\
\midrule
$131072$ (all 4s) & 32{,}781 & 1{,}929{,}679 & 11{,}999 & 1{,}941{,}678 \\
full chain (all 4s) & 65{,}533 & 3{,}521{,}967 & 11{,}999 & 3{,}533{,}966 \\
maximum score (2s) & 129{,}333 & 658{,}984 & 147 & 659{,}131 \\
\bottomrule
\end{tabular}
\end{center}

Measured runs match: the browser reports $1.91$--$1.97$ million undos for
the $131072$ line. The maximum-score line is four times longer but three
times cheaper in undos, because a 2 is nine times as likely as a 4.

\section{Honest play}\label{sec:honest}

For the full board we implemented an expectimax player (\textsc{Genius})
in the tradition of Xiao's \cite{nneonneo2014}: rows and columns are
table lookups over ranks, the heuristic scores empty cells, available
merges, monotonicity and a tile-sum penalty, with a pull toward the chosen
corner; depth adapts to the number of distinct tiles. Against Richardson's
\cite{richardson2014} Evil generator, which places each new tile on the
edge the last move packed against beside the largest neighbours, the
search models that rule exactly instead of averaging over luck. We ported
Richardson's four players (Smart, Algorithm, Priority, Random) and the
generator to the same array engine for a like-for-like comparison. Ten
games each, bottom-right corner:

\begin{center}\small
\begin{tabular}{@{}lrrr@{}}
\toprule
player & regular tiles & Evil tiles & regular, undo to escape death \\
\midrule
\textsc{Genius} & 16384 $\times1$, 8192 $\times5$, 4096 $\times3$, 1024 $\times1$; 130{,}131 & 4096 $\times6$, 2048 $\times4$; 56{,}850 & 32768 $\times2$ (2 games) \\
Smart & 4096 $\times3$, 2048 $\times6$, 1024 $\times1$; 46{,}029 & 1024 $\times4$, 512 $\times4$, 256 $\times2$; 8{,}410 & 16384 $\times3$, 8192 $\times2$ \\
Algorithm & 512 $\times5$, 256 $\times2$, 128 $\times3$; 4{,}256 & 256/128/64; 2{,}011 & 1024 $\times6$, 512 $\times4$ \\
Priority & 512 $\times1$, 256 $\times6$, 128 $\times2$, 16 $\times1$; 2{,}952 & 128/64; 1{,}284 & --- \\
Random & 128 $\times6$, 64 $\times3$, 32 $\times1$; 983 & 128/64/32; 508 & 256 $\times9$, 512 $\times1$ \\
\bottomrule
\end{tabular}
\end{center}

Semicolons separate the largest-tile histogram from the average score. The
last column is the undo button used as a human uses it: only to escape
game over, each death unwinding a doubling number of moves, giving up
after forty deaths without a new best score.

\section{Open problems}

\begin{enumerate}[nosep]
\item Conjecture~\ref{conj}: a $4\times4$ play with sixteen spawned 4s.
\item The maximum score of $4\times3$ (predicted $180{,}180$) and
$2\times5$, within reach of layered enumeration.
\item Whether Theorem~\ref{thm:score} is tight on \emph{every} board, and
whether a uniform construction exists.
\item Minimum undos: the expected undo count above is for the shipped
lines; the line minimising expected undos (which would prefer spawns with
many empty cells) is a different optimisation.
\end{enumerate}

\section*{Reproducibility}

All code is in the repository \cite{superintelligence2026}. \texttt{node
test/solve.js 3x3} reproduces Table~\ref{tab:exact} and
Table~\ref{tab:honest} ($3\times3$ takes about five minutes);
\texttt{PERFECT=1 node test/run.js} replays each shipped line through the
game engine and asserts the exact counts (\texttt{GOAL=spiral|score},
\texttt{ORIENT=col}); \texttt{node test/gen\_perfect.js} regenerates a line
from nothing; \texttt{node test/honest.js} and \texttt{node
test/honest\_drive.js} produce the honest-play table.

\bibliographystyle{plain}
\bibliography{refs}

\end{document}
